\section{The split rows and the global contribution from the columns}

% BEGIN G81S1
\paragraph{G81S1. The two finite ranks on a conjugate pair.}
Use the hub of F78D2 at a twisted split middle-range index $i$ as in
M81B1, with $p=\ell_i$, $\pi\in Z_{\rm good}(i)$. Write
$W=\rho_i(\pi)/\chi_1$ and $W^*=W^\vee(1)$. The declaration
pairs each column $u$ with the conjugate row $w$. Indeed $p\ge5$,
$p$ avoids $q,d,g_0,g_1$, and $-d$ is a nonzero square modulo $p$ by
the split class. Its two roots lift to distinct roots of $X^2+d$
in $\mathbb Q_p$ by Hensel's lemma. The host CM type is distinguished
by the value on $\sqrt{-d}$. At a fixed completion all embeddings
have the same value on this element, since its image is already in
$\mathbb Q_p$. Complex conjugation exchanges the two roots, hence
exchanges the two types of places and preserves their local degrees.
Selecting the type prescribed by $\mathrm{epJ}$ for columns, and
the opposite type for rows, gives a finite disjoint covering and
satisfies both clauses of $\mathrm{hypDecl}$. Thus the witnessed
case of $\mathrm{dcOf}$ applies. Every witness has the same members
by those clauses. S80J1 supplies $\mathrm{ConclusionJa}$, and
the sign selected from it has its own labelled jump identities.
Use those identities for this very sign in the lists of S80J1.

At the column all embeddings have the upward list, and the rank
block has the cloud and the $m$ lows. Its cloud jumps are positive.
At the row every embedding has the downward list; its cloud is
strictly below $-1$, so the cloud of the dual $W^*$ there is
positive. The low weights of $W$ lie outside $[-1,t]$ in either
orientation by the two gaps in S80J1. Thus $0$ is absent from
$W$ at the column and from $W^*$ at the row. Both representations
are crystalline: the sum at the selected point is crystalline
away from the type primes by S80J3, and the rank summand, its
dual, and its Tate twist are crystalline. A fixed rational
vector would give a trivial subrepresentation and hence a jump
at zero by strictness. For either representation the dimension
of its de Rham quotient by $\mathrm{Fil}^0$ is the sum of all
negative jump multiplicities over the embeddings. Only the $m$
lows per embedding can contribute. The crystalline dimension
formula in J80H1 therefore gives
\[
       j_W(u)\le d_um,\qquad
       \dim H^1_f(E_w,W^*)\le d_um,\qquad
       d_u=[E_u:\mathbb Q_p]=[E_w:\mathbb Q_p].
\]
The inequalities hold over every tested coefficient extension.
% END G81S1

% BEGIN G81S2
\paragraph{G81S2. The split row length at the stored exponent.}
Fix the row $w$ conjugate to $u$, any good coefficient field with
integers $O'$, and any of the native surjections onto a
matrix-identified lattice $T$ for $W^*$. Write $k_{O'}$ for
$\operatorname{length}_{O'}(O'/p)$ and $a=X.e_S(i)$.
Every jump of $W^*$ has absolute value at most
$h=\lfloor2b\rfloor+1$ by M81B1 and the dual rule $z\mapsto-1-z$.
At each embedding there are at most $r$ distinct jumps, and none
equals zero. On the unramified field $E_w$, W79H1 and W79H5 give
the de Rham comparison and the good level-one Wach lattice
required by A79L5--A79L8. Apply that invariant theorem to $T$,
which is allowed to be any of the tested lattices. The calculation
of S79P5 uses only $p,h,r$, the absence of zero and the unramified
degree bound. Thus its per-embedding exponent is at most
$2+r(4+2\lfloor\log_p(2h)\rfloor)\le X.c_0(i,u)$, since
the canonical $c_0$ has the same formula at every place. Writing
$V=T[1/p]$, this kills $(V/T)^{G_{E_w}}$ by $p^{X.c_0(i,u)}$.
This module embeds in $p^{-X.c_0(i,u)}T/T$, so its length is at
most $r X.c_0(i,u) k_{O'}$. This also proves the original
$\mathrm{H0ColumnLengthBound}$ for the conjugate row lattice.

Set $P=H^1_f(E_w,T)$, the preimage of the rational finite
subspace. Its torsion is a quotient of $(V/T)^{G_{E_w}}$ by
the invariants sequence for $T\to V\to V/T$. Its free rank
is at most $d_um$ by G81S1: finite generation and
rationalization follow from the integral argument of A81R1.
L81S4 applies the dual-row lifting condition on the whole
family at the prescribed $a$. In particular the image of
finite source cohomology in $P$ contains $p^aP$, since the
target and source of a tested row are permitted subquotients
of the dual sums (the target is a direct summand rationally).
The quotient by that image has length at most
$a d_um k_{O'}+\operatorname{length}_{O'}P_{\rm tors}$,
by the decomposition of a finite module over the DVR into its
free and torsion parts. The cocycle definition of row length
also quotients by integral coboundaries and therefore gives
exactly this cohomological cokernel. Its normalized length is
at most $d_um a+rX.c_0(i,u)$. This proves
$\mathrm{ColengthRowPlaceBound}$ and, assigning this bound to
each conjugate row, $\mathrm{LengthBoundAtPlaces}$ with
$\mathrm{lenRowSplit}$ on the actual split declaration.
% END G81S2

% BEGIN G81S3
\paragraph{G81S3. A regular character corner when roots coincide modulo the prime.}
On this same split index, replace the family matrix at $\pi$ by
the direct sum of its given integral rank and character matrices.
Its rational equivariant isomorphism with the original point
preserves characteristic polynomials, point evaluations and the
Hecke algebra. Elsewhere keep the original matrices. Local tests
on subobjects of rational sums transport by this conjugation.
For row and column lifting the tested lattices in those subobjects
are arbitrary and the standard target rank lattice is still the
original one. The Hecke evaluation ideal and its support are
unchanged by this normalization.

Use D79G1--D79G3 to construct the integral orthogonal projectors
and the rank-$t$ character corner $e_B$ throughout the family.
Since $E_w$ is unramified, inertia surjects onto the full
cyclotomic Galois group over $\mathbb Q_p$: each finite layer
is totally ramified, and compactness gives the inverse limit.
Take $g_0$ with cyclotomic value $1+p$. Its $t$ character
values at $\pi$ are distinct in characteristic zero. Set
$X_0=e_BR(g_0)e_B$. We can perturb this corner integrally away
from $\pi$ to have simple spectrum on its image at every $x$.

To prove this, fix $x\ne\pi$ and extend to an algebraic closure
of the coefficient field. Every absolutely irreducible global
constituent of $R_x$ in a composition series occurs at most once.
Indeed such a finite series, with its constituent isomorphisms,
is defined over a finite coefficient extension: choose bases of
the subspaces and write down their finitely many entries and the
entries of any isomorphism in the algebraic closure. An invariant
subspace over this field is invariant before passing to its
closure by the same matrix equations. Restrict to a fixed place
over $p$ and use the crystalline comparison of S80J3. Exactness
and strictness of the filtered periods give additivity of every
jump multiplicity on the series. A nonzero factor has a jump at
every embedding. Repeating an isomorphic factor would repeat that
jump, contrary to the simple spectrum of Hodge weights in S80J1
(twisting all factors by the same character merely shifts it).

Let $A_x$ be the span algebra of the group matrices over the
algebraic closure. Its semisimple quotient acts on the simple
constituents as the product of their full matrix algebras. This
follows from the structure theorem for finite-dimensional
algebras over an algebraically closed field; every simple
factor is seen on the faithful module, by passage to a composition
series and nilpotence on the resulting filtration. Thus the corner
$e_B A_x e_B$ surjects onto the corresponding product of matrix
corners, by multiplying a lift on both sides by $e_B$. Choose
diagonal matrices with all entries distinct across the nonzero
corner spaces. The filtration of the image of $e_B$ on the
original module has precisely these corner spaces as its successive
quotients, each once. A lift has distinct eigenvalues on that
image. It follows that the discriminant of the characteristic
polynomial there is a nonzero polynomial on the corner.

The Hecke algebra contains a function supported at $x$ with
nonzero value there, by the separating-products proof of G79E1.
After multiplication by this function the integral corners
$e_B R(g)e_B$ still span the entire corner at $x$ over the
algebraic closure. Take finitely many that span. Consequently the
discriminant of $X_0$ plus a linear combination of these supported
corners is a nonzero polynomial in the coefficients. An infinite
integral domain is Zariski dense in the affine space over its
algebraic closure of fractions: induct on the number of variables
using the finite root bound in one variable. Choose coefficients
in the original DVR avoiding zero. Repeat separately for all
$x\ne\pi$, retaining the initial value at $\pi$.

Enlarge the coefficient field once to split all the resulting
polynomials. The eigenvalues are integral because the polynomials
are monic and integral. D79G5 constructs the ordered-root cover
and its free monomial basis over the extended Hecke algebra:
distinctness in the fraction field is enough for injectivity of
evaluation there, and the ordered splitting algebra is torsion-free.
D79G6 produces the integral first-character polynomial $E_1$
of rank one at every point. Write $\delta$ for its product of
root differences evaluated at the distinguished lift of $\pi$.
The ratios at that lift are $(1+p)^{-j}$, $1\le j<t$.
Since $p$ is odd, their differences from one have normalized
valuations $1+v_p(j)$, by the binomial formula for powers of
$1+p$. Hence $v_p(\delta)=v_\delta(i)$ in the native convention
$s_0(i)=1$. All statements here concern the extended coefficient
ring if a splitting field was needed.

For any block write $C_b$ for the global cocycle module constructed
from these matrices as in G79E7. The determinant proof
G79E11--G79E12 gives the evaluation length bounded by the sum
of the block lengths. It only needs rank one off support, the
free ordered-root basis and surjectivity onto the separate blocks,
now supplied above. The off-support assertion of D79G7 persists
after coefficient extension by scalar extension of the common
linear equations for a fixed vector. The extended evaluation
ideal is the old ideal extended to the new DVR: restrictions
to the complement form a finite matrix, and flat extension
preserves its kernel and the image of evaluation at $\pi$.

For the intersection with global coboundaries, repeat the equation
of G79E7 before division by $\delta$. On any block other than the
first character, intertwining gives
$\delta\sum_j b_j(\pi')z(g_j)=0$. For a coboundary
$z(g)=(M(g)-1)v$ the sum equals $-\delta v$, so $\delta^2 v=0$.
In block dimension $d_b$ the coboundaries in $C_b$ consequently
have length at most $2d_b v_\delta(i)k_{O'}$; they are images of
a submodule of the $\delta^2$-torsion of $(L'/O')^{d_b}$.
The first character has zero coboundaries. The total normalized
loss in passing from cocycles to global classes is therefore at
most $2(r+t-1)v_\delta(i)$.
% END G81S3

% BEGIN G81S4
\paragraph{G81S4. The precise remaining column estimate for the split inequality.}
Keep the coefficient ring $O'$ and the global block classes of
G81S3. For the rank block take the kernel of restriction to the
local quotients by $H_f^1$ at all rows. For each character take
the kernel of restriction to all auxiliary inertias in $Q_i$
and to the local quotients by $H_f^1$ at all rows. Denote the
resulting modules by $K_b$. Define $D_{\rm col}$ to be the
image of $\bigoplus_b K_b$ under simultaneous restriction in
each block to the quotients of local cohomology by $H_f^1$ at
all columns. Its normalized length is denoted by $d_{\rm col}$,
allowing infinite length. Write $c$ for the number of columns,
$a=X.e_S(i)$, $v=v_\delta(i)$, and $h_0$ for the source's maximum
normalized invariant length of the nontrivial character ratios
on the columns. A sufficient remaining estimate on these very
modules is
\[
\begin{split}
 d_{\rm col}\le{}&
 \sum_{u\in\mathrm{col}}
        \{a+d_um(a+v)+2v+h_W(i,u,\pi)\}\\
 &+tc h_0+\{4n_V(c+1)-2(r+t-1)\}v .
\end{split}                                                     \tag{*}
\]
This condition involves the common Hecke column modules after the
row and auxiliary restrictions. Neither a global lattice-index
inequality nor an estimate over independently chosen representations
is assumed in the implication below.

To prove the implication, I81G3--I81G4 apply at each row $w$
to any finite list of the rank cuts from G81S3. Their graph
lattice uses the same point complement. The cut presentation by
the columns of $E_1$ in G79E13 uses only their integrality and
intertwining. There is no division by $\delta$ there. G81S2
therefore bounds the normalized loss on this row by
$d_um a+rX.c_0(i,u)$ for its conjugate column $u$.
The class before imposing row finiteness on each character has
auxiliary inertia loss at most $\sum_{P\in Q_i}z_P$ as in G79E8.
At a row apply I81G5 with M81B1--M81B3 to its local cut
presentation. Its normalized loss is at most $(d_w+1)X.c_B(i)$.
This includes primes with $p-1\mid j+1$ on the $j$th block,
since the extra integral $H^2$ term is cyclic as proved in M81B3.
The joint restriction loss is bounded by the sum of the individual
bounds. The kernel after the further column restrictions has
normalized length at most
\[
                    s_\pi+v_p(h_E)+\sum_{0<j<t}u_j .
\]
Indeed the rank classes now lie in the defining Selmer group
as in I81G6. The first character is everywhere unramified;
for the others M81B3 makes the finite local class zero above $p$.
D79G7 and G79E8 ensure inertia vanishes at the other prescribed
places, and the trace to the class group and the cyclotomic
Tate--Shafarevich comparison of G79E6 apply. Here $u_j$ denotes
the native length for the inverse character of exponent $j$.
These are already modules of cohomology classes by G81S3, so a
global coboundary contributes no further loss.

All normalized original terms survive the common coefficient
enlargement. As in I81G2, a finite free basis of $O'$ over the
original $O$ writes the continuous cochains as finitely many
copies before extension; it respects restriction, invariants,
inertia, coboundaries and the crystalline finite parts. Hence
it computes the same base change for the Selmer kernel.
The factor multiplying finite lengths is the relative ramification
index, also the factor multiplying the length of $O/p$.
For positive twists the restriction-of-scalars proof in G79E6
gives exactly their stated $u_j$. The same base change applies
to $z_P$ and to the evaluation ideal by G81S3. Killing and odd
finiteness at the chosen good point supply the required finite
terms and the normalized total auxiliary bound $p^{-2}$.

Add the block lengths using G81S3. Insert (*) and the row bounds,
and replace $d_w+1$ by $d_w+2$ in the nonnegative character
errors. The global coboundary loss uses precisely the allowance
subtracted in (*). For each column the terms involving $a$
become $(2d_um+1)a$ after adding its conjugate row. The remaining
pair terms are $d_um v+2v+h_W(i,u,\pi)+rX.c_0(i,u)$.
The other errors are
\[
    t X.c_B(i)\sum_{w\in\mathrm{row}}(d_w+2)
            +4 n_V(c+1)v+tc h_0 .
\]
By the canonical split identities in DFD8 and N78T2 the sum
of the $(2d_um+1)a$ terms is the exact $\mathrm{sumOver}$ of
$X.\mathrm{len0}(i,\pi)$, and all the remaining errors are
$X.\mathrm{errorBoundHub}(i,\pi)$. This proves the original split
$\mathrm{StarAD}$ inequality at the chosen point under (*).
The rows, the ordinary and dual integral lifting clauses, and
the character condition at each row are supplied above. Bounding
the actual simultaneous image in (*) by the displayed allowance
remains open in this argument.
% END G81S4
